\clearpage
\VerdierPhysicalPageAnchor{0142}%
\begingroup
\small
\setlength{\parskip}{0pt}%
\setlength{\emergencystretch}{1.5em}%

\VerdierSemanticAnchor{ch2-2-2-11}
\noindent\textbf{Corollaire 2.2.11.}\quad
\begingroup\itshape
\VerdierSemanticAnchor{ch2-2-2-11-setup}%
Soient $\mathcal D$ une catégorie triangulée et
$\mathcal B\subset\mathcal D$ une sous-catégorie triangulée pleine.

\VerdierSemanticAnchor{ch2-2-2-11-a}
\par\noindent a) Le noyau du foncteur de passage au quotient
$Q\colon\mathcal D\longrightarrow\mathcal D/\mathcal B$ est la plus
petite sous-catégorie triangulée strictement pleine saturée qui contienne
$\mathcal B$.

\VerdierSemanticAnchor{ch2-2-2-11-b-order-isomorphism}
\par\noindent b) L'application
$\mathcal B\longmapsto S(\mathcal B)$ qui à toute sous-catégorie
triangulée strictement pleine saturée associe son système multiplicatif
saturé (2.1.8) est un isomorphisme d'ensembles ordonnés de l'ensemble des
sous-catégories triangulées strictement pleines saturées, ordonné par
inclusion, dans l'ensemble des systèmes multiplicatifs saturés de
$\mathcal D$, ordonné par inclusion.

\VerdierSemanticAnchor{ch2-2-2-11-b-inverse}
L'application inverse associe à tout système multiplicatif saturé le noyau
du foncteur de localisation correspondant.

\VerdierSemanticAnchor{ch2-2-2-11-c-factorization}
\par\noindent c) Un foncteur exact (resp. cohomologique)
$F\colon\mathcal D\longrightarrow\mathcal D'$ dont le noyau contient la
sous-catégorie $\mathcal B$ se factorise d'une manière unique par
$Q\colon\mathcal D\longrightarrow\mathcal D/\mathcal B$ en un foncteur
$G\colon\mathcal D/\mathcal B\longrightarrow\mathcal D'$ qui est exact
(resp. cohomologique).

\VerdierSemanticAnchor{ch2-2-2-11-c-morphism-lifting}
Un morphisme de foncteurs exacts (resp. cohomologiques) dont les noyaux
contiennent la sous-catégorie $\mathcal B$ provient d'une manière unique
d'un morphisme entre les factorisations.
\endgroup

\VerdierSemanticAnchor{ch2-2-2-11-easy-consequences}
Ce sont des conséquences faciles du théorème (2.2.6).

\VerdierSemanticAnchor{ch2-2-2-12}
\smallskip
\noindent\textbf{2.2.12. Commentaires sur la propriété TRIV.}
\VerdierSemanticAnchor{ch2-2-2-12-pretriangulated-category}%
Appelons \emph{catégorie pré-triangulée} une $\mathbb Z$-catégorie stricte
munie d'une famille de triangles possédant les propriétés TRI, TRII et
TRIII.

\VerdierSemanticAnchor{ch2-2-2-12-induced-notions}
On définit de la même façon que pour les catégories triangulées les notions
de \emph{foncteurs exacts}, \emph{foncteurs cohomologiques},
\emph{systèmes multiplicatifs}.

\VerdierSemanticAnchor{ch2-2-2-12-localization}
Soit $S$ un système multiplicatif (2.1.4) d'une catégorie pré-triangulée
$\mathcal D$. La catégorie localisée $\mathcal D(S^{-1})$ admet une et une
seule structure de catégorie pré-triangulée telle que le foncteur de
localisation soit exact.

\VerdierSemanticAnchor{ch2-2-2-12-kernel-subcategory}
La sous-catégorie pleine des objets de $\mathcal D$ qui sont annulés par le
foncteur de localisation est une sous-catégorie pré-triangulée pleine.

\VerdierSemanticAnchor{ch2-2-2-12-triv-necessity}
On a besoin dans la pratique de savoir, inversement, associer à une
sous-catégorie pré-triangulée un système multiplicatif saturé convenable.
Le rédacteur ne sait le faire qu'en utilisant la propriété TRIV.

\VerdierSemanticAnchor{ch2-2-2-12-triv-role}
Cette propriété, de nature assez technique, a donc surtout servi, pour
l'instant, à obtenir le corollaire (2.2.11) à partir du théorème (2.2.6).

\VerdierSemanticAnchor{ch2-2-3}
\medskip
\noindent\textbf{2.3. Propriétés du foncteur de localisation.}

\VerdierSemanticAnchor{ch2-2-3-1}
\smallskip
\noindent\textbf{Proposition 2.3.1.}\quad
\itshape Soient $\mathcal D$ une catégorie triangulée,
$\mathcal B$ une sous-catégorie triangulée pleine.\upshape

\VerdierSemanticAnchor{ch2-2-3-1-p143-boundary}
\endgroup

